\clearpage\thispagestyle{empty}
\pdfbookmark[0]{Quellen zum Kometenbericht vom 10. April 1826}{biela1826-sources}
\begin{center}\textbf{Quellen zum Kometenbericht vom 10. April 1826}\end{center}
\noindent Diese KI-Übertragung ins heutige Deutsch folgt dem Bericht \textit{Comet} in \href{https://gdz.sub.uni-goettingen.de/id/PPN236007467}{den \textit{Werken}, Band~VI} (Göttingen, 1874), S.~451--453. Der Bericht erschien ohne Unterschrift unter \textit{Göttingen} in den \textit{Göttingischen gelehrten Anzeigen}, Nr.~57 (10. April 1826), S.~\href{https://api.digitale-sammlungen.de/iiif/image/v2/bsb10538703_00567/full/1550,/0/default.jpg}{561}, \href{https://api.digitale-sammlungen.de/iiif/image/v2/bsb10538703_00568/full/1550,/0/default.jpg}{562}, \href{https://api.digitale-sammlungen.de/iiif/image/v2/bsb10538703_00569/full/1550,/0/default.jpg}{563}, \href{https://api.digitale-sammlungen.de/iiif/image/v2/bsb10538703_00570/full/1550,/0/default.jpg}{564} und \href{https://api.digitale-sammlungen.de/iiif/image/v2/bsb10538703_00571/full/1550,/0/default.jpg}{565}. Die Links führen zum Exemplar der Bayerischen Staatsbibliothek. Der Bericht verweist auf die früheren Anzeigen in \hyperlink{native-422}{Nr.~28} und \hyperlink{native-425}{Nr.~83} von 1819.
\par\smallskip\noindent Der eingeklammerte Verweis auf S.~417 und 420 gehört zur Werkausgabe. Die Tabellen folgen dieser Ausgabe. Dazu gehören die erste Längenangabe ohne Sekundenzeichen und das Kürzel O ohne Punkt in der vierten Längenzeile. Sowohl in der Zeitschrift als auch in der Werkausgabe stehen an den betreffenden Stellen $4\frac{3}{4}$ und $3\frac{3}{8}$ Jahre. KI-Übertragung und KI-Satz.
\clearpage
